% !TEX root = chapter-standalone.tex


\section{Associativity}
\label{sec:associativity}

In the previous section we discussed how the notion of interface is central to series compositions of many different kinds. All of the examples we saw also share a further common feature, which is known as \emph{associativity} (or \emph{the associative property}). This is a generalization of the notion of associativity that we are already familiar with from number systems, semigroups, monoids, etc.

Let us explain associativity of series composition using the electrical cables example. Instead of focusing only on the series composition of two electrical cables, suppose we wish to compose three compatible cables together:
\begin{equation}
    \label{eq:three-cables}
    \cord{\TypeIrish}{a}{\TypeSwiss}
    \qquad
    \cord{\TypeSwiss}{b}{\TypeUSA}
    \qquad
    \cord{\TypeUSA}{c}{\TypeChinese}
\end{equation}

If we connect two composable cables at a time, there are two steps needed and two possible variants of how we do it, depending on the order of our actions. On the one hand, we could connect the first two cables together, and then connect their composite to the last cable:
\begin{equation}
    \label{eq:cables-assoc-left}
    \prfperiod{
        \prftree{
            \cord{\TypeIrish}{a}{\TypeSwiss}
        }{
            \cord{\TypeSwiss}{b}{\TypeUSA}
        }{
            \cord{\TypeIrish}{a+b}{\TypeUSA}
        }
    }{
        \cord{\TypeUSA}{c}{\TypeChinese}
    }{
        \cord{\TypeIrish}{(a+b)+c}{\TypeChinese}
    }
\end{equation}

Or, on the other hand, we could first connect the last two cables together, and then connect the first cable to their composite:
\begin{equation}
    \label{eq:cables-assoc-right}
    \prfperiod{
        \cord{\TypeIrish}{a}{\TypeSwiss}
    }{
        \prftree{
            \cord{\TypeSwiss}{b}{\TypeUSA}
        }{
            \cord{\TypeUSA}{c}{\TypeChinese}
        }{
            \cord{\TypeSwiss}{b+c}{\TypeChinese}
        }
    }{
        \cord{\TypeIrish}{a+(b+c)}{\TypeChinese}
    }
\end{equation}

Either procedure results in the same final outcome: a cable with the same types of ends, and with the same length, since $(a+b)+c = a+(b+c)$. 

\begin{equation}
    \label{eq:cables-assoc-all}
    \prfperiod{
        \cord{\TypeIrish}{a}{\TypeSwiss}
    }{
        \cord{\TypeSwiss}{b}{\TypeUSA}
    }{
        \cord{\TypeUSA}{c}{\TypeChinese}
    }{
        \cord{\TypeIrish}{a+b+c}{\TypeChinese}
    }
\end{equation}

That is the essence of associativity: there might be different possible orders in which we implement a nested sequence of pairwise series compositions, but the overall result will be the same, no matter the route taken. This is analogous to the associativity that we are used to from, say, the addition of numbers; the order in which we add up numbers pairwise doesn't matter for the end result:
\begin{equation}
    (\ela + \elb) + \elc = \ela + (\elb + \elc).
\end{equation}

If we are summing up just three numbers, then there are only two different possible routes to summing them up using pairwise sums. If we have many more numbers, then there will be more different possible routes, but associativity still guarantees that the final result from all of them will be the same. For instance, for five numbers we have
\begin{align*}
      & \ (((\ela + \elb) + \elc) + \eld) + \ele     \\
    = & \ ((\ela + \elb) + \elc) + (\eld + \ele)     \\
    = & \ (\ela + \elb) + (\elc + (\eld + \ele))     \\
    = & \ \ela + (\elb + (\elc + (\eld + \ele))).
\end{align*}

The same is true for the associativity of series composition: for any finite number of electrical cables that are composable in series, the order in which we compose them pairwise into one long cable doesn't matter for the end result.

This pattern also applies to all of the examples of series composition discussed in the previous section. We highlight a few cases, to illustrate. In each case, we consider three things that are composable in series, and we compare the two ways of composing them.

\paragraph{Sets and functions}
Let $\setA$, $\setB$, $\setC$, and $\setD$ be sets, and consider functions $\mapa \colon \setA \sto \setB$, $\mapb \colon \setB \sto \setC$, and $\mapc \colon \setC \sto \setD$. Both ways of composing them give functions from $\setA$ to $\setD$, and for every $\ela \setin \setA$ we have
\begin{equation}
    \begin{aligned}
        ((\mapa \mthen \mapb) \mthen \mapc)(\ela) & = \mapc((\mapa \mthen \mapb)(\ela)) = \mapc(\mapb(\mapa(\ela))), \\
        (\mapa \mthen (\mapb \mthen \mapc))(\ela) & = (\mapb \mthen \mapc)(\mapa(\ela)) = \mapc(\mapb(\mapa(\ela))).
    \end{aligned}
\end{equation}
Since the two functions have the same value at every element of $\setA$, they are equal:
\begin{equation}
    \label{eq:assoc-functions}
    \prfperiod{
        \mapa \colon \setA \sto \setB
    }{
        \mapb \colon \setB \sto \setC
    }{
        \mapc \colon \setC \sto \setD
    }{
        (\mapa \mthen \mapb) \mthen \mapc = \mapa \mthen (\mapb \mthen \mapc)
    }
\end{equation}
In both cases, we first apply $\mapa$, then $\mapb$, and then $\mapc$; the parentheses only record which two of the functions we composed first.

\paragraph{Real coordinate spaces and real matrices}
Consider matrices $\matmor{A} \colon \reals^k \mto \reals^m$, $\matmor{B} \colon \reals^m \mto \reals^n$, and $\matmor{C} \colon \reals^n \mto \reals^p$. Recall from \cref{sec:interfaces} that the series composition $\matmor{A} \mthen \matmor{B}$ is given by the matrix product $\matmor{B}\matmor{A}$. The two ways of composing the three matrices in series therefore give the $p \times k$ matrices
\begin{equation}
    (\matmor{A} \mthen \matmor{B}) \mthen \matmor{C} = \matmor{C}(\matmor{B}\matmor{A})
    \qqand
    \matmor{A} \mthen (\matmor{B} \mthen \matmor{C}) = (\matmor{C}\matmor{B})\matmor{A}.
\end{equation}
These are equal, because matrix multiplication is associative. Writing $\matmor{A}_{rl}$ for the entry of $\matmor{A}$ in row $r$ and column $l$, and similarly for the other matrices, both products have the entry
\begin{equation}
    \sum_{j=1}^{n} \sum_{r=1}^{m} \matmor{C}_{ij} \, \matmor{B}_{jr} \, \matmor{A}_{rl}
\end{equation}
in row $i$ and column $l$. Thus:
\begin{equation}
    \label{eq:assoc-matrices}
    \prfperiod{
        \matmor{A} \colon \reals^k \mto \reals^m
    }{
        \matmor{B} \colon \reals^m \mto \reals^n
    }{
        \matmor{C} \colon \reals^n \mto \reals^p
    }{
        (\matmor{A} \mthen \matmor{B}) \mthen \matmor{C} = \matmor{A} \mthen (\matmor{B} \mthen \matmor{C})
    }
\end{equation}
This is the matrix counterpart of the associativity of function composition: both matrices encode the linear function which transforms a vector $\vectob{v} \setin \reals^k$ into $\matmor{C}(\matmor{B}(\matmor{A}\vectob{v}))$, that is, which first applies $\matmor{A}$, then $\matmor{B}$, and then $\matmor{C}$.

\paragraph{Components in co-design}
Let us extend the example of the motor and the battery from \cref{sec:interfaces} by a third component: a drivetrain (the gears and shafts that transmit the rotation of the motor to the wheels), which provides a certain torque at the wheels, and which requires a certain torque from the motor. There are two ways to obtain the overall system, which provides torque at the wheels and requires mass. We could first combine the drivetrain and the motor into a subsystem which provides torque at the wheels and requires electric current, and then connect this subsystem to the battery. Or we could first combine the motor and the battery into a subsystem which provides torque of the motor and requires mass, and then connect the drivetrain to this subsystem. Both ways result in the same design problem. In general, for posets $\posA$, $\posB$, $\posC$, and $\posD$, we have (\cref{lem:dp-series-laws})
\begin{equation}
    \label{eq:assoc-dp}
    \prfperiod{
        \adpa \colon \funposA \profto \resposB
    }{
        \adpb \colon \funposB \profto \resposC
    }{
        \adpc \colon \funposC \profto \resposD
    }{
        (\adpa \dpthen \adpb) \dpthen \adpc = \adpa \dpthen (\adpb \dpthen \adpc)
    }
\end{equation}
Intuitively, in both cases a functionality in $\posA$ is feasible with a requirement in $\posD$ exactly when we can find intermediate resources in $\posB$ and in $\posC$ (in the example: a torque of the motor and an electric current) with which each of the three components is feasible.

\paragraph{Journeys}
Let us extend our journey from Zürich to Milan by train ($\mora$) and from Milan to Florence by car ($\morb$) with a third leg: from Florence to Rome by train ($\morc$). We could first combine the first two journeys into a journey from Zürich to Florence, and then continue to Rome. Or we could first combine the last two journeys into a journey from Milan to Rome, and then precede it by the train journey from Zürich to Milan. Either way, we travel the same three legs in the same order, so we obtain the same journey from Zürich to Rome. In general, for locations $\Obja$, $\Objb$, $\Objc$, and $\Objd$:
\begin{equation}
    \label{eq:assoc-journeys}
    \prfperiod{
        \mora \colon \Obja \mto \Objb
    }{
        \morb \colon \Objb \mto \Objc
    }{
        \morc \colon \Objc \mto \Objd
    }{
        (\mora \mthen \morb) \mthen \morc = \mora \mthen (\morb \mthen \morc)
    }
\end{equation}
If we think of a journey as a list of legs, each of which ends where the next one begins, then series composition of journeys is concatenation of lists, and the associativity of composing journeys is the associativity of concatenating lists.

\paragraph{The general pattern}
In all of these examples, we have encountered the same equation: whenever we have three things that are composable in series,
\begin{equation}
    \label{eq:assoc-generic}
    \prfperiod{
        \mora \colon \Obja \mto \Objb
    }{
        \morb \colon \Objb \mto \Objc
    }{
        \morc \colon \Objc \mto \Objd
    }{
        (\mora \mthen \morb) \mthen \morc = \mora \mthen (\morb \mthen \morc)
    }
\end{equation}
As with sums of numbers, it follows that for any finite sequence of things composable in series, all the different ways of composing them pairwise give the same result. This is why we may omit the parentheses, and simply write $\mora \mthen \morb \mthen \morc$. Associativity will be one of the conditions in the definition of a category (\cref{def:categorymain}).
